\documentclass[11pt]{article}
\usepackage[T1]{fontenc}
\usepackage{lmodern}
\usepackage[margin=1in]{geometry}
\usepackage{amsmath,amssymb,amsthm,mathtools,booktabs,array,longtable}
\usepackage{microtype,xcolor,xurl,enumitem}
\usepackage[colorlinks=true,linkcolor=blue!45!black,urlcolor=blue!55!black,citecolor=blue!45!black]{hyperref}
\usepackage{fancyhdr}
\pagestyle{fancy}\fancyhf{}\lhead{A smaller nine-variable consistency equation}\rhead{23 October 2026}\cfoot{\thepage}
\setlength{\headheight}{14pt}
\setlength{\emergencystretch}{2em}
\setlist{nosep,leftmargin=*}
\newcommand{\N}{\mathbb N}
\newcommand{\Z}{\mathbb Z}
\newcommand{\Pos}{\mathbb Z_{>0}}
\newcommand{\Con}{\operatorname{Con}}
\newcommand{\pop}{\operatorname{pop}}
\newcommand{\Gcap}{G_{\mathrm{cap}}}
\newcommand{\Sc}{S_{\mathrm c}}
\newcommand{\Tc}{T_{\mathrm c}}
\newcommand{\Jc}{J_{\mathrm c}}
\newcommand{\Mc}{M_{\mathrm c}}
\newcommand{\Dc}{D_{\mathrm c}}
\newcommand{\Ac}{A_{\mathrm c}}
\newtheorem{theorem}{Theorem}[section]
\newtheorem{lemma}[theorem]{Lemma}
\newtheorem{proposition}[theorem]{Proposition}
\newtheorem{corollary}[theorem]{Corollary}
\theoremstyle{remark}\newtheorem{remark}[theorem]{Remark}
\title{\textbf{Free Pell coordinates and a smaller\newline nine-variable equation for ZFC consistency}\\[0.6em]
\large Exact degree $4\,686\,812$; 560 binary arithmetic operations}
\author{Research note prepared for Vladimir Reshetnikov and ProveIt\\
\small AI-assisted mathematical construction and independent computational audits}
\date{23 October 2026}

\begin{document}
\maketitle
\begin{abstract}
We construct a fully specified integer-coefficient polynomial in nine
strictly positive integer unknowns whose zero-existence condition is
equivalent to the repository's fixed 31-variable refutation system.
Composing with its stated ordinary compiler and K215 calibration theorems
gives the desired ZFC-provable equivalence between absence of zeros and
$\Con(\mathrm{ZFC})$. The new polynomial has exact total degree
$4\,686\,812$ and a division-free evaluator with 560 live binary operations.
Against the lowest-degree nine-variable candidate found in the incoming
archives, $7\,344\,328$ in 738 operations, the reductions are
$36.18\%$ in degree and $24.12\%$ in arithmetic size.

The principal new idea is to replace a degree-$\deg R$ Pell coordinate by
a free positive coordinate. A positive-margin condition already implied
by the final norm reconstructs its old lower bound. A small quotient
translation makes that reasoning valid before the margin is known to be
positive. A second change removes an unnecessary leading binary block
from the carry code. A shared evaluator for twenty of the twenty-two
quadratic rows supplies most of the operation saving. We also give
fifteen-positive-variable equations with 73 operations at degree 54 and
74 operations at degree 46.

All new transformations are proved as ordinary mathematics. Exact
coefficient data, literal circuits, source bindings, and reproducible
finite checks accompany the paper. The result is not a claim of global
minimality or a new end-to-end Lean/Rocq certificate; the inherited
compiler dependency and the distinction from an object-language ZFC
derivation are kept explicit.
\end{abstract}
\newpage
\tableofcontents
\newpage

\section{Result, comparison, and meaning of ``simpler''}

There is no known single ordering that captures simplicity of Diophantine
equations. Variable count, total degree, arithmetic circuit size, coefficient
size, and the length of a correctness proof can move in different directions.
This note improves the degree and arithmetic size together at the
nine-positive-variable point. It does not claim to reduce the established
nine-natural-variable existence bound to eight, nor to give an optimal
polynomial among all representations.

The comparison is frozen at ProveIt commit
\texttt{4b4349064a750ba6347293b55874a7f92abdd0db}.
Both the placed collection and all eight incoming ZIP archives were inspected.
This matters: several incoming improvements are absent from the active short
circuit documentation. In particular, the correct short-operation comparator
is 66, not 67. Table~\ref{tab:frontier} distinguishes those dimensions.

\begin{table}[ht]
\centering\small
\begin{tabular}{@{}llrr@{}}
\toprule
Construction & Witness domain & Exact degree & Operations\\
\midrule
Placed package 45 & $\Pos^9$ & $129\,875\,980$ & 788\\
Incoming selective quadratic & $\Pos^9$ & $7\,344\,328$ & 738\\
Incoming sparse Pell, fastest listed & $\Pos^9$ & $15\,502\,748$ & 548\\
\textbf{This note, primary} & $\boldsymbol{\Pos^9}$ & $\boldsymbol{4\,686\,812}$ & \textbf{560}\\
This note, original-$H$ alternative & $\Pos^9$ & $4\,686\,812$ & 562\\
This note, nonnegative translate & $\N^9$ & $4\,686\,812$ & 567\\
\midrule
Incoming squared quotient & $\Pos^{15}$ & 157 & 66\\
Incoming squared quotient SOS & $\Pos^{15}$ & 80 & 72\\
\textbf{This note, $Q_{54}$} & $\boldsymbol{\Pos^{15}}$ & \textbf{54} & \textbf{73}\\
\textbf{This note, $Q_{46}$} & $\boldsymbol{\Pos^{15}}$ & \textbf{46} & \textbf{74}\\
Incoming gap minimum & $\Pos^{15}$ & 44 & 75\\
\bottomrule
\end{tabular}
\caption{Audited comparison scope. The nine-variable primary improves the
lowest-degree incoming point in both columns. The 548-operation incoming
point remains a different tradeoff. No unrestricted-integer assertion is
made for the positive-domain polynomials.}
\label{tab:frontier}
\end{table}

Let $\mathsf{Ref}_{\mathrm{ZFC}}$ mean that there is a finite syntactic
derivation of falsity from the repository's ZFC axiom predicate, and let
$\Con(\mathrm{ZFC})$ be its negation. Our main, entirely local mathematical
result is a constructive equivalence between the new polynomial and the
specified finite source system. The interpretation of that source system
as $\mathsf{Ref}_{\mathrm{ZFC}}$ is the named inherited compiler theorem,
not an inference from sampled evaluations.

\begin{theorem}[Construction and endpoint]\label{thm:main}
Let $p_1,\ldots,p_{22}$ be the fixed quadratic polynomials of
Appendix~\ref{app:source}, with the complete package-31 coefficient tuple
specified in Section~\ref{sec:coefficients}. The polynomial $\mathcal P$
defined below belongs to $\Z[b,g,h,s,w,C,i,j,n]$ and satisfies
\[
 \exists\,(b,g,h,s,w,C,i,j,n)\in\Pos^9:\ \mathcal P=0
 \quad\Longleftrightarrow\quad
 \exists\,\mathbf x\in\N^{31}:\ p_1(\mathbf x)=\cdots=p_{22}(\mathbf x)=0.
\]
It has exact total degree $4\,686\,812$ and the supplied primary circuit
has 332 multiplications, 173 additions, and 55 subtractions.
Using the repository's ordinary source-compiler and K215 calibration
theorems, the equivalence is formalizable in ZFC and yields
\[
 \boxed{\quad
 \mathrm{ZFC}\ \vdash\
 \bigl[\neg\exists\,\mathbf u\in\Pos^9:\mathcal P(\mathbf u)=0\bigr]
 \ \longleftrightarrow\ \Con(\mathrm{ZFC}).\quad}
\]
\end{theorem}

The nine variables in this statement are the only quantified unknowns.
Every other symbol introduced below is a polynomial abbreviation or a
fixed integer. Fixed powers are expanded by multiplication in the literal
circuit. The theorem does not assume consistency and does not decide
whether the equation has a zero.

\section{The explicit equation}\label{sec:equation}

\subsection{Fixed source data and carry polynomials}

Put
\[
 \kappa=2\,134\,940,\qquad
 \lambda=2^\kappa,\qquad
 \chi=2^{16\,012\,039},\qquad B=\lambda b^2.
\]
The fixed integer-coefficient polynomial $\Dc$ and the two fixed lists
$w_1,\ldots,w_{31}$ and $t_1,\ldots,t_{22}$ are defined in
Section~\ref{sec:coefficients}. Set
\begin{align}
 \ell&=\sum_{r=1}^{31}B^{w_r},&
 \Ac&=\sum_{r=1}^{22}B^{t_r},&N_0&=B^{184},\label{eq:source1}\\
 \Mc&=N_0-1-(b-1)\ell,&
 \Jc&=2(1+g)^2\Dc(B)+\Ac,\label{eq:source2}\\
 \Sc&=g+N_0\Jc,&
 \Tc&=\Mc+N_0(B-2)\Ac,\label{eq:source3}\\
 L&=bB^{5741},&
 R&=(\Sc+\Tc)L+L-1-\Sc,\label{eq:source4}\\
 N&=\chi B^{2966},&
 \Gcap&=32bB^{183}.\label{eq:source5}
\end{align}
Here $L$ is the binary block length and $N$ is the threshold supplied to
Jones's ratio theorem. They are different polynomials. In particular,
the old leading $L^2$ term is absent from $R$.

\subsection{Pell registers and the free coordinate}

The coordinate $C$ in the following display is a free positive unknown.
It is not defined to be $2R+1+bw+\phi$.
\begin{align}
 U&=N^2w,&Y&=N^2s,&M_0&=RY,\notag\\
 P&=2M_0^2U,&K&=R+1+h(P-1),&A&=M_0(U+1),\notag\\
 B_1&=2R+1,&W&=bw,&J&=4A-5,\notag\\
 Q&=3WC-2(W^2-1),&D&=(A^2-1)C^2+1,\notag\\
 E&=iJDC^2,&F&=(A^2-1)E^2+1,\notag\\
 G&=A+F(CD+1-A),&
 H&=C+F\bigl(B_1+(j-1)C\bigr),\notag\\
 I&=(G^2-1)H^2+1.\label{eq:pell}
\end{align}
The use of $j-1$ is intentional: the CRT-absorbed parameter ranges over
$\N$, while the displayed nine-variable presentation uses only positive
inputs. This shift is charged in the 560-operation ledger.

Define
\begin{align}
 a_0&=(P^2-1)K^2+1,&a_1&=64DFI,\notag\\
 m&=\bigl[K^2-4(C-KY)^2\bigr]\Gcap-8gK^2,\notag\\
 k_*&=iJC,&T_*&=k_*(8D^2-4D-1),\notag\\
 V_0&=PK,&V_1&=T_*GH,\notag\\
 \Omega&=Jm,&
 \Theta&=J(n+W^2)+Q-\Omega(V_0+V_1).\label{eq:final-registers}
\end{align}
The requested equation is
\begin{equation}\label{eq:main}
 \boxed{\qquad
 \mathcal P
  =\bigl[\Theta^2-\Omega^2(a_0+a_1)\bigr]^2
       -4\Omega^4a_0a_1=0.
 \qquad}
\end{equation}
Equations~\eqref{eq:source1}--\eqref{eq:final-registers} are definitions,
not additional equations with independent unknowns. Thus
\eqref{eq:main} is one ordinary polynomial equation in nine unknowns.
The fixed coefficient recipe contains no consistency-dependent selection,
unbounded proof search, variable exponent, or unspecified r.e. index.

\subsection{A slightly longer form with a simpler congruence interface}

For the 562-operation alternative, replace $H$ in~\eqref{eq:pell} by
$H=B_1+jC$ and put
\[
 u=JF,\quad v=QF+J(H-C),\quad
 \Omega=um,\quad \Theta=u(n+W^2)+v-\Omega(V_0+V_1).
\]
Keep the same norm~\eqref{eq:main}. Every input then appears directly as
a positive coordinate. The degree is unchanged. This version is useful
for checking the precise connection to the original Jones ratio theorem;
the two presentations are existence-equivalent, not claimed to be the
same polynomial on the same tuples.

\section{Coefficient binding and finite selective code}\label{sec:coefficients}

\subsection{A coherent, materialized tuple}

The four source ports are $\beta,K_0,M_C,M_F$. We use the entire
\texttt{wang943-mixed-121x8801} tuple from placed package 31, including
the already shifted $M_F$ port. The archive contains all four compressed
little-endian magnitude payloads, their hashes, the original binding
manifest, and the coefficient-construction proof. Their bit lengths are
\[
 1\,064\,921,\quad1\,070\,004,\quad1\,064\,921,\quad1\,064\,922.
\]
In particular $\beta=2^{1\,064\,921}-1$. The other three fixed integers
are given by those finite payloads and their reproducible native compiler
recipe. An expanded decimal printing is unnecessary for a complete
definition and would obscure the construction. The supplied data do not
leave any port to be chosen later.

Let $p_r(\mathbf x)$ be the residual in row $r$ of Appendix~\ref{app:source}
after every one of its positive source coordinates is replaced by $x_i+1$.
For $|\alpha|\leq2$, define
\[
 c_{r,\alpha}=(2-|\alpha|)!\prod_i\alpha_i!\,
                   [\mathbf x^\alpha]p_r(\mathbf x).
\]
There are 136 nonzero row-monomial records, using 71 distinct monomials.
Each $c_{r,\alpha}$ is an explicitly supplied polynomial in the four fixed
ports. Bounding a recipe $\sum_e a_e\prod_s q_s^{e_s}$ by
$\sum_e|a_e|2^{\sum_s e_s\operatorname{bitlength}(q_s)}$ gives
\[
 |c_{r,\alpha}|<M_{\rm coef}=2^{2\,134\,927},
 \qquad \lambda=8192M_{\rm coef}.
\]
The finite verifier reconstructs the shifted rows and these recipes from
the literal source, rather than trusting the declared counts.

\subsection{The selected weights and row positions}

In the variable order of Appendix~\ref{app:source}, use
\begin{align*}
 (w_1,\ldots,w_{31})={}&(9,25,23,139,35,151,168,21,91,181,94,52,150,180,7,135,\\
 &38,88,183,106,178,136,20,138,8,64,105,154,37,137,120).
\end{align*}
Their maximum is 183. Assign targets to rows as follows:
\begin{center}\small
\begin{tabular}{@{}rrrrrrrrrrr@{}}
\toprule
Row &16&2&9&14&21&7&15&22&13&3\\
Target &70&86&451&685&1050&1296&1530&1874&2212&2358\\
\midrule
Row &19&20&1&10&17&8&4&6&11&18\\
Target &2696&2780&2792&3157&3290&3634&3929&4238&4384&4527\\
\midrule
Row &12&5&&&&&&&&\\
Target &4865&5190&&&&&&&&\\
\bottomrule
\end{tabular}
\end{center}
Write $t_r$ for the target assigned to row $r$, and
$e(\alpha)=\sum_i\alpha_iw_i$. Define the fixed polynomial
\begin{equation}\label{eq:fixedD}
 \Dc(z)=\sum_{r,\alpha}c_{r,\alpha}z^{t_r-e(\alpha)}.
\end{equation}
Every exponent is nonnegative. The leading exponent is 5190 and its
coefficient is exactly 2.

The certificate checks, for every supported pair $(r,\alpha)$ and every
multi-index $\gamma$ of degree at most two,
\begin{equation}\label{eq:isolation}
 t_r-e(\alpha)+e(\gamma)=t_s
 \quad\Longleftrightarrow\quad r=s\ \text{and}\ \alpha=\gamma.
\end{equation}
There are $136\cdot528=71\,808$ interactions. This is a complete finite
check of a stated combinatorial condition, not a probabilistic test.
It implies the coefficient identity
\begin{equation}\label{eq:extract}
 [z^{t_r}]\left(1+\sum_i x_i z^{w_i}\right)^2\Dc(z)
       =2p_r(\mathbf x).
\end{equation}
Indeed the coefficient of a padded monomial in the square is
$2!/((2-|\alpha|)!\prod_i\alpha_i!)$, which cancels the factorial weight.
The selective code is inherited from the incoming selective-quadratic
archive; its reuse is not claimed as a new combinatorial discovery.

\section{A smaller carry code: remove the leading block}\label{sec:carry}

\subsection{Bounds valid before any equation holds}

For all positive $b$, $B\geq\lambda$, so $B-1>M_{\rm coef}$ and $B>128$.
The positive leading coefficient of $\Dc$ dominates all its lower signed
coefficients:
\[
 B^{5190}<\Dc(B)<\frac{M_{\rm coef}B^{5191}}{B-1}.
\]
The digits of $\Mc$ are $B-b$ at the 31 selected weights and $B-1$
elsewhere in its 184-digit range. Thus $0\leq\Mc<N_0$,
$\Jc,\Sc>0$, and $\Tc\geq0$. In particular
\begin{equation}\label{eq:Rpositive}
 R=(L-1)\Sc+L(\Tc+1)-1\geq L-1\geq\max(8,b).
\end{equation}
Also $N\geq\max(8,b)$ and $\Gcap>0$. These unconditional facts are what
the later ratio argument needs. No inequality $R\geq N$ is assumed.

On the branch $g<4bB^{183}$, integrality gives $1+g\leq4bB^{183}$,
and $g<N_0$ because $B>4b$. Put $Z=B^{5557}$. Then
\[
 2(1+g)^2\Dc(B)
 <\frac{32M_{\rm coef}b^2 B^{5557}}{B-1}<\frac Z{128},
 \qquad \Ac<\frac Z{128},
\]
so $0<\Jc<Z/64$. Consequently
\[
 \Sc<N_0(Z/64+1)<B^{5741}/32.
\]
Since $(B-2)\Ac<B^{5191}$ and $B^{366}>64$,
\[
 \Tc<N_0(1+B^{5191})<B^{5741}/32.
\]
It follows that
\begin{equation}\label{eq:blockbound}
 0\leq\Sc,\Tc,\qquad\Sc+\Tc<L=bB^{5741},
\end{equation}
including the boundary $b=1$.

\subsection{The carry test detects precisely the source zeros}

Suppose $b$ is a power of two; then so is $B$. The low binary mask says
\[
 g\mathbin{\&}\Mc=0
 \quad\Longleftrightarrow\quad
 g=\sum_i x_iB^{w_i},\qquad0\leq x_i<b.
\]
The coefficient at the padding position zero is fixed at one by
$(1+g)^2$; zero is not one of the variable weights.

For such $g$, put
$h_t=[z^t](1+\sum_i x_i z^{w_i})^2\Dc(z)$. The coefficient bound gives
\[
 |h_t|<M_{\rm coef}(1+\textstyle\sum_i x_i)^2
          <1024M_{\rm coef}b^2=B/8.
\]
Before ordinary radix normalization, the digits of $\Jc$ are
$2h_t+\delta_t$, where $\delta_t=1$ at a target and zero elsewhere.
An incoming borrow is always $0$ or $-1$. At a target, a zero $h_t$
therefore leaves the normalized digit $0$ or $1$. A positive $h_t$ leaves
a digit at least 2 and below $B/2$; a negative $h_t$ leaves a digit above
$B/2$. The mask $B-2$ accepts exactly the first case. This proves
\[
 \Jc\mathbin{\&}(B-2)\Ac=0
 \quad\Longleftrightarrow\quad h_{t_r}=0\text{ for all }r.
\]
The borrow assertion follows directly by induction from
$|h_t|<B/8$; it does not use a bit operation on negative integers.
The evaluated $\Jc$ is already positive and below $Z$.

The low and high blocks in~\eqref{eq:source3} do not overlap. Using
\eqref{eq:extract}, we obtain
\begin{equation}\label{eq:carrysource}
 \Sc\mathbin{\&}\Tc=0
 \quad\Longleftrightarrow\quad
 g\text{ encodes a common nonnegative zero of all }p_r.
\end{equation}
Conversely, for a source zero choose any sufficiently large power of two
$b$ and use its radix encoding $g$. The source has $p_5(\mathbf0)=1$,
so that encoding is positive. Moreover $g<bB^{183}$ and therefore
$32g<\Gcap$, the stronger bound needed for the completeness direction of
Jones's theorem.

\subsection{Exact two-block valuation and its parity}

\begin{lemma}[Two blocks without a leading sentinel]\label{lem:twoblock}
Let $L=2^u$ and $S,T\geq0$ with $S+T<L$. Set
\[
 R=(S+T)L+(L-1-S).
\]
Then
\[
 v_2\binom{2R}{R}
 =u+\pop(T)-\bigl(\pop(S)+\pop(T)-\pop(S+T)\bigr).
\]
In particular $2^{u+\pop(T)}\mid\binom{2R}{R}$ if and only if $S\&T=0$.
\end{lemma}
\begin{proof}
The two binary blocks do not overlap, and the low block is the $u$-bit
complement of $S$. Thus $\pop(R)=\pop(S+T)+u-\pop(S)$.
Legendre's identity $v_2(r!)=r-\pop(r)$ gives
$v_2\binom{2R}{R}=\pop(R)$. Finally
$\pop(S)+\pop(T)-\pop(S+T)$ is nonnegative and vanishes exactly when
the addition has no carries. This proves both assertions.
\end{proof}

Write $b=2^k$ and $B=2^v$, so $v=\kappa+2k$. The low mask has
$184v-31k$ set bits, and the 22 high target digits each have $v-1$ set
bits. Therefore
\[
 \pop(\Tc)=206v-31k-22,\qquad
 u=\log_2 L=k+5741v.
\]
Their sum is
\begin{equation}\label{eq:threshold}
 u+\pop(\Tc)=5947\kappa+11864k-22
       =2\log_2\bigl(\chi B^{2966}\bigr).
\end{equation}
The exponent of $\chi$ is $15\kappa/2-11=16\,012\,039$.
Together, Lemma~\ref{lem:twoblock} and~\eqref{eq:carrysource} show that
the source is represented by the three numerical requirements
\[
 b\text{ is a power of two},\qquad
 N^2\mid\binom{2R}{R},\qquad g<4bB^{183}.
\]
The incoming comparator used $2bB^{5741}$ and an additional $L^2$ term.
Deleting that term and halving the block size changes the threshold by
two binary places, explaining the factor-two change in $\chi$.
The proof above checks this exact change; it is not a heuristic degree
estimate. At $b=1$ the low mask would force $g=0$, so the enlarged positive
parameter domain adds no zero.

\section{The free-coordinate theorem}\label{sec:freeC}

\subsection{The numerical interface imported from Jones}

We use the numerical ratio theorems, not a black-box application of the
hard-wired universal-polynomial formula. In the repository their precise
interfaces are \texttt{Jones1982.Nine.ratio\_sound} and
\texttt{ratio\_complete} in \texttt{NineRatio.lean}, corresponding to the
Pell and binomial analysis of Jones~\cite{Jones1982}. With
$R,N\geq8$, $1\leq b\leq\min(R,N)$ and $\Gcap>0$, the standard registers
use
\[
 C=B_1+W+\phi,\quad H_{\rm old}=B_1+j_{\rm old}C,
 \qquad\phi\geq0,\quad j_{\rm old}>0.
\]
Their four conditions are
\begin{equation}\label{eq:ratio-conditions}
 a_0=\square,\qquad DFI=\square,\qquad
 JF\mid QF+J(H_{\rm old}-C),\qquad m>0.
\end{equation}
The soundness theorem implies a power-of-two $b$, the central-binomial
divisibility, and $8g<\Gcap$. The completeness theorem supplies witnesses
from the same power and divisibility conditions when $32g<\Gcap$.
These asymmetric bounds are precisely those established in the preceding
section. All source data satisfy the unconditional size hypotheses.

The next proof explains why the lower-bound slack $\phi$ can be removed
from the actual polynomial. Recovering it after using the ratio theorem
would be circular. We instead recover it from the final norm before that
theorem is invoked.

\subsection{Unconditional positive majorants}

Assume only the allowed coordinate domains and the unconditional source
bounds. Then $U,Y\geq64$, $K\geq P+R>1$, and $A\geq2$. Moreover
\[
 A=RY(U+1)>U=N^2w\geq bw=W.
\]
Because $C\geq1$ is a positive independent input,
\[
 D\geq A^2>W^2,\qquad D>C^2,\qquad WC<D.
\]
The last inequality follows by squaring or from $W,C<\sqrt D$.
Also $CD+1-A>0$, so $G>A$, and $H>1$ in either presentation.
Consequently $a_0,a_1$ are positive on every allowed assignment.

The first strict majorant is immediate:
\begin{equation}\label{eq:maj0}
 V_0^2-a_0=K^2-1>0.
\end{equation}
Since $E=k_*DC$ and $(A^2-1)C^2=D-1$, we have
$F=(D-1)k_*^2D^2+1$. The following exact identity supplies the second:
\begin{align}
 T_*^2-64DF&=(8D+1)k_*^2-64D\notag\\
           &=8D(k_*^2-8)+k_*^2>0.\label{eq:sharp}
\end{align}
Indeed $k_*=iJC\geq3$. Thus
\begin{equation}\label{eq:delta}
 \Delta_1:=V_1^2-a_1
 =\bigl[(8D+1)k_*^2-64D\bigr](GH)^2
                         +64DF(H^2-1)>0.
\end{equation}
The positive signs of $V_0,V_1$ matter: a squared inequality by itself
would not justify the later positivity argument. The majorant
\eqref{eq:sharp} is inherited from the incoming sparse-Pell work.

\subsection{A quotient shift makes the free positive \texorpdfstring{$C$}{C} safe}

The free coordinate may make $Q$ negative. The offset $W^2$, already
computed inside $Q$, repairs that problem with one extra addition:
\begin{equation}\label{eq:positive-numerator}
 J(n+W^2)+Q
   =Jn+(J-2)W^2+3WC+2>0.
\end{equation}
For completeness we also need an upper bound on the translated quotient.
Because $W\geq1$ and $J\geq3$,
\[
 Q/J\leq WC<D,\qquad W^2+Q/J<2D.
\]
The elementary inequality
$8D^2-4D-1\geq3D$, together with $k_*\geq3$, gives
$T_*\geq9D$. Hence $V_1>2D$ and
\begin{equation}\label{eq:quotient-bound}
 W^2+Q/J<V_1.
\end{equation}
No positivity of $Q$ was assumed. In the original-$H$ version, use
$u=JF$ and $v=QF+J(H-C)$. Then $H-C>0$, the translated numerator is
still positive, and
\[
 W^2+v/u=W^2+Q/J+(H-C)/F<2D+H<V_1.
\]

\subsection{Soundness: prove the margin first}

Suppose $\mathcal P=0$. If $m=0$, then $\Omega=0$ and
$\mathcal P=[J(n+W^2)+Q]^4>0$, a contradiction. For $m\ne0$, the
factorization of the quartic norm over the reals provides signed real
roots $r_0^2=a_0$, $r_1^2=a_1$ with
\begin{equation}\label{eq:root-eq}
 n+W^2+Q/J=mZ_*,\qquad Z_*=V_0+V_1+r_0+r_1.
\end{equation}
The strict positive majorants imply $Z_*>0$ even before either root is
known to be rational. The left side is positive by
\eqref{eq:positive-numerator}; therefore $m>0$.

Since $g>0$ and $\Gcap>0$, the definition of $m$ now yields
\[
 4(C-KY)^2<K^2,\qquad C>K(Y-\tfrac12)>2K.
\]
Because $K\geq P+R$ and $P>W$, this implies
\begin{equation}\label{eq:restorephi}
 C>2(P+R)>2R+1+W=B_1+W.
\end{equation}
Thus the missing old coordinate $\phi=C-B_1-W$ is a positive integer.
This is the structural saving: the polynomial did not have to build that
lower bound into the expression for $C$.

To justify the radical merger, we must also exclude cancellation of two
identical nonsquare radicals. We have
\[
 P<A^2,\qquad D>C^2,\qquad F\geq A^2,\qquad I\geq G^2\geq A^2.
\]
Together with $C>K$, these imply
\begin{equation}\label{eq:separate}
 DFI>A^4C^2>P^2C^2>P^2K^2>a_0,
 \qquad a_1=64DFI>a_0.
\end{equation}
Therefore $x=r_0+r_1=\Theta/\Omega$ is a nonzero rational number.
The identity $2xr_0=x^2+a_0-a_1$ makes $r_0$ rational, and then so is
$r_1$. A rational number whose square is an integer is an integer: reduce
its numerator and denominator to coprime integers. Hence both roots are
integral, and $r_1/8$ is likewise an integer square root of $DFI$.
Equation~\eqref{eq:root-eq} now forces $J\mid Q$.

It remains to restore the old $j$ parameter. Put $j_0=j-1\geq0$.
The definitions give $F\equiv1\pmod C$ and $F\equiv1\pmod J$. Define
\[
 j_{\rm old}=1+\frac{F-1}{C}B_1+Fj_0>0.
\]
This is an integer and gives $H=B_1+j_{\rm old}C$ and $F\mid H-C$.
Since $\gcd(F,J)=1$, the recovered $J\mid Q$ is exactly what is needed
for the combined divisibility in~\eqref{eq:ratio-conditions}. The source
sizes, restored $\phi$, square conditions, and positive margin now permit
the Jones soundness theorem. It gives the carry conditions, which give
a common source zero by Section~\ref{sec:carry}.

\subsection{Completeness: reconstruct the final quotient}

Start from a common source zero. Choose a sufficiently large power-of-two
$b$, encode $g$, and apply Jones completeness using $32g<\Gcap$.
Retain the old $C=B_1+W+\phi$ as the new positive coordinate.
The old combined divisibility, and $\gcd(F,J)=1$, give $J\mid Q$ and
$F\mid H_{\rm old}-C$. Let $z=(H_{\rm old}-C)/F$. Then $z>0$,
$z\equiv B_1\pmod C$, and $0<B_1<C$; consequently
$z=B_1+j_0C$ for a unique $j_0\geq0$. Take $j=j_0+1$.
This gives exactly the new formula for $H$, without changing the radicands.

Choose positive integral roots $r_0,r_1$ of $a_0,a_1$, and set
\[
 n=m(V_0+V_1+r_0+r_1)-W^2-Q/J.
\]
All quantities on the right are integers. By~\eqref{eq:quotient-bound}
and $m\geq1$, the right side is strictly positive. The corresponding
linear factor of the norm vanishes. This completes the proof of the
source-system equivalence in Theorem~\ref{thm:main}.

\begin{remark}[What the witness map asserts]
Completeness chooses a new final quotient using positive square roots.
It does not assert that subtracting $W^2$ from every old norm witness
preserves positivity. Similarly, the CRT parameter is allowed to be zero
internally; its positive-coordinate shift cannot be silently deleted.
\end{remark}

\section{Exact degree and arithmetic size}\label{sec:metrics}

\subsection{Resolving the leading cancellation}

The most important degrees of the primary, CRT-absorbed presentation are
listed in Table~\ref{tab:degrees}. They refer to the nine unknowns only;
the four source ports and the fixed scales have degree zero.
\begin{table}[ht]
\centering\small
\begin{tabular}{@{}lr@{\qquad}lr@{\qquad}lr@{}}
\toprule
Register & Degree & Register & Degree & Register & Degree\\
\midrule
$\Dc(B)$ &10380 &$R$ &22233 &$N$ &5932\\
$A$ &45963 &$C$ &1 &$D$ &91928\\
$F$ &367714 &$G$ &459643 &$H$ &389947\\
$a_0$ &320246 &$a_1$ &2158822 &$m$ &184221\\
$V_0$ &160123 &$V_1$ &1079411 &$\Omega$ &230184\\
$\Delta_1$ &1883038 &$\text{norm inner}$ &2343406 &$\mathcal P$ &4686812\\
\bottomrule
\end{tabular}
\caption{Exact degrees certified by a full multivariate upper bound and
an explicitly nonzero diagonal leading coefficient.}
\label{tab:degrees}
\end{table}

A naive maximum-and-sum degree propagation through~\eqref{eq:main}
would overestimate its degree. Let
$E_0=J(n+W^2)+Q$. The exact identity
\begin{align}
 \Theta^2-\Omega^2(a_0+a_1)
 ={}&E_0^2-2E_0\Omega(V_0+V_1)\notag\\
 &+\Omega^2\bigl[(K^2-1)+\Delta_1+2V_0V_1\bigr]\label{eq:degreeidentity}
\end{align}
uses~\eqref{eq:maj0} and~\eqref{eq:delta}. In the primary construction,
$\Omega^2\Delta_1$ is the unique term of greatest degree on the right,
namely $2\,343\,406$. Its square has degree $4\,686\,812$; the remaining
term $4\Omega^4a_0a_1$ has strictly smaller degree. This gives a valid
multivariate upper bound and isolates the leading form to be checked.

On the diagonal $b=g=h=s=w=C=i=j=n=t$, the source leaders are
\[
 [R]_{22233}=\lambda^{11115}(\lambda+4),\qquad
 [N]_{5932}=\chi\lambda^{2966}.
\]
The two terms $\Sc$ and $\Tc$ have the same degree and positive leaders;
this is why $\lambda+4$ appears after deleting the old leading block.
Substituting into the positive defect in~\eqref{eq:delta} and then
\eqref{eq:degreeidentity} gives the exact diagonal leading coefficient
\begin{equation}\label{eq:leader}
 \boxed{\quad
 2^{94}\lambda^{2\,343\,092}(\lambda+4)^{104}\chi^{400}>0.
 \quad}
\end{equation}
Thus the upper bound is attained for the actual fixed coefficients.
This is an exact nonvanishing argument, not interpolation or sampling.
The independent degree-certificate circuit uses the expanded defect
identities, so unresolved top-degree cancellations cause verification to
fail instead of being guessed away.

For a useful parametric check, let $r=\deg R$, $v=\deg N$, and let $W_c$
be the maximum source weight. For either $H$ presentation, the resulting
degree in this range is
\[
 104r+400v+8W_c+316.
\]
Here $(r,v,W_c)=(22233,5932,183)$ gives the announced result.
Keeping the original leading block would give $4\,763\,044$; restoring
the old composite $C$ as well gives $6\,692\,108$. These are checks on
the contribution of each new step, not separate simplicity claims.

\subsection{Shared evaluation of twenty quadratic rows}

The degree improvements do not by themselves explain the 560-operation
evaluator. For arithmetic size, reuse the common 60-operation quadratic
source schedule to evaluate many of the rows in~\eqref{eq:fixedD}
together. Put $W_c=183$, $u_0=B^{W_c}$, and $u_i=B^{W_c-w_i}$.
Evaluate each original positive source coordinate as $u_i+u_0$, aligning
lower-degree additions by powers of $u_0$. For row $r$ this computes the
degree-two homogenization of $p_r$ at $(u_0,u_1,\ldots,u_{31})$.

The factorial weight is two on a diagonal monomial, including the
constant padding square, and one on a mixed monomial. Adding one further
copy of each diagonal term therefore gives the required weighted row.
For a target $t_r\geq2W_c$, multiplication by $B^{t_r-2W_c}$ produces
that row's contribution to $\Dc(B)$ with no negative powers.

Exactly twenty targets meet this condition. The two low targets, 70
and 86, are evaluated directly from their sparse monomials. They cannot
be forced through the homogeneous schedule using a negative exponent.
This mixed evaluator preserves the unchanged selective code and gains
sharing without moving any target. The independent checker expands all
22 source rows exactly, including the transport-coordinate adapter between
the shared schedule and the authoritative original source.

The primary circuit has 332 multiplications, 173 additions, and 55
subtractions, for 560 operations. Every operation is live. Multiplication
by a fixed numeral, every squaring, and every multiplication used to build
a fixed power is charged. Reuse of an existing register is free. Fixed
integer coefficients are input leaves, as in the repository's comparison
metric; producing their binary payload is a separate finite preprocessing
task. The certificate circuit is deliberately longer and is not used for
the operation-count claim.

\section{Proof status, reproducibility, and limits}\label{sec:status}

\subsection{What is proved here and what is imported}

The free-coordinate and quotient-shift lemmas, the removal of the leading
binary block, and the local short-equation synthesis are proved in this
paper. The fixed selective code and sharpened Pell majorant are inherited
from incoming submissions, with their identities checked again. The
Jones ratio theorem is a separate numerical result with a repository Lean
formalization. The interpretation of the fixed 31-variable system uses
the repository's ordinary native-window compiler, coherent tuple adapter,
and K215-to-syntactic-refutation calibration.

Those layers should not be conflated. The archive supplies no new
end-to-end kernel proof of the 560-operation equation, and the finite
audits cannot establish the universal compiler theorem. It also does not
export a derivation inside the repository's object-language ZFC proof
calculus. The stated ZFC-provable corollary is an ordinary mathematical
composition of the identified arithmetic and compiler theorems. The
active repository separately contains a closed, checked nine-variable
construction; its proof status does not automatically certify this much
smaller explicit circuit.

The distinction is relevant to incorporation: the new relative
source-system equivalence can be formalized independently of the native
compiler. Its hypotheses are just the explicit numerical source bounds
and the existing Jones ratio interfaces. No new consistency assumption,
oracle, nonconstructive choice of coefficients, or additional mathematical
axiom is part of the construction.

\subsection{Reproducible finite evidence}

The package is self-contained for its numerical construction and audits;
Python's standard library suffices. From the extracted archive root run
\begin{verbatim}
python3 verify_all.py
python3 -O verify_all.py
\end{verbatim}
The driver checks the source payloads and declared fixed coefficient
recipes, the full selective-isolation condition, the homogenized row
identities, every operand and live gate, the Pell and norm identities,
and the exact leading-degree certificate. The short-equation audit
expands its complete leading homogeneous forms. Additional modular
comparisons exercise separately compiled evaluators and are reported as
transcription checks, not as universal proofs. Finite binary fixtures
also check the two-block valuation identity, including its smallest
boundaries.

The main files are the literal positive circuit, its independently
arranged certificate circuit and degree ledger, the coefficient binding,
the preserved source data and proofs, and the two short-equation circuits.
The top-level README gives exact paths and evaluation commands. The
hash manifest fixes the delivered bytes. The source pin is recorded so
later repository changes cannot silently redefine the comparison.

To rebuild this paper from the archive root:
\begin{verbatim}
cd report
pdflatex -interaction=nonstopmode -halt-on-error report.tex
pdflatex -interaction=nonstopmode -halt-on-error report.tex
\end{verbatim}

\subsection{Boundaries of the conclusion}

These equations concern positive or explicitly translated nonnegative
integer solutions. Substituting arbitrary integers for all inputs is not
licensed by the proofs. A conversion to unrestricted integers would
require a separate domain encoding and a new cost ledger.

The fixed coefficients remain very large, even though their compressed
recipes and payloads are finite and reproducible. A 560-operation
arithmetic circuit is not a claim of fast evaluation on large arbitrary
inputs: bit complexity is a different measure. Finally, the note does
not prove consistency, inconsistency, or an absolute minimum degree.
Under the usual hypotheses of the second incompleteness theorem, a
consistent ZFC cannot prove the displayed absence of zeros, since the
equivalence would then prove its own consistency. No two-sided
independence statement is inferred from consistency alone.

\appendix
\section{Additional short equations: 73 operations at degree 54}
\label{sec:short-synthesis}

The incoming archives contain useful improvements that are absent from
the active P76 documentation: a 66-operation equation of degree 157,
a 72-operation equation of degree 80, and an ordinary gap equation with
75 operations at degree 44. We retain these as the comparison baseline.
The following synthesis adds the two tradeoffs
\[
 (\text{operations},\text{degree})=(73,54),\qquad (74,46),
\]
each using fifteen strictly positive integer variables. It combines the
incoming minimum-norm argument and squared-modulus reconstruction, and
regroups the main norm to save one further operation. The finite verifier
checks every literal arithmetic gate and exact leading homogeneous form.
The semantic compiler remains an identified inherited mathematical
dependency; no new proof-assistant or object-language ZFC certificate is
asserted in this supplement.

\subsection{The exact polynomial}

Use the complete package-31 coefficient tuple, with $B=2^{1064921}$,
$\beta=B-1$, and its already shifted $M_F$. The tuple and all four binary
payloads are included in \texttt{short/coefficients31/}.
The fifteen positive unknowns are
\[
 J,F,\alpha,z,f,h,i,T,s,w,\tau,\eta,\zeta,e,\gamma.
\]
Define polynomial abbreviations, with no additional unknowns:
\begin{align*}
 q&=\beta J+1,& C&=q-F-\alpha,\\
 X&=wq,&Y&=sq^3,&E&=XY,\\
 k&=\eta+\zeta,&c&=kY+\eta,&a&=E+Y,\\
 H&=4a+3,&\Delta&=a^2+H,&D&=X+ac+\gamma H,\\
 L&=E(kY),&
 R&=((q-1)(q-F)+\alpha)(q^2-1)+(M_C+qM_F)J,\\
 V&=f^2(cT-R)-c,&U&=\Delta(f^2-1),&I&=ic^2.
\end{align*}
The five factors and strong residual are
\begin{align*}
 N_0&=\tau^2-L(L+k),&N_1&=D^2-\Delta c^2,\\
 N_a&=(V+e)^2-Ue(2V+e),&W_s&=I^2-U,\\
 N_k&=k-hE-R,&N_t&=(K_0+w)C+(q-F)-z(q-1).
\end{align*}
The main short equation is
\begin{equation}\label{eq:short54}
 \boxed{Q_{54}=N_0^2+N_a^2+W_s^2+(N_1N_kN_t-1)^2-2=0.}
\end{equation}
The lower-degree alternative is
\begin{equation}\label{eq:short46}
 \boxed{Q_{46}=N_0^2+N_1^2+N_a^2+W_s^2+(N_kN_t-1)^2-3=0.}
\end{equation}
The subscripts denote exact degree. Their literal evaluators have 73
and 74 operations, respectively. The principal evaluator has
39 multiplications, 19 additions and 15 subtractions; the alternative
has one more addition. Multiplication by a fixed numeral is charged.

\subsection{The minimum argument and complete local equivalence}

The following facts hold on every positive assignment, before imposing
either equation:
\begin{equation}\label{eq:min-blocks}
 N_0^2\geq1,\qquad N_a^2+W_s^2\geq1.
\end{equation}
Equality cases are $N_0=1$ and $(N_a,W_s)=(1,0)$, respectively.
For the first assertion, write $n=XY^2\geq2$. The discriminant $n(n+1)$
is nonsquare, so $N_0\ne0$. A hypothetical solution of
$x^2-n(n+1)t^2=-1$ has
$nt<|x|<(n+\tfrac12)t$. Multiplication by
$2n+1-2\sqrt{n(n+1)}$ produces a new positive second coordinate
$(2n+1)t-2|x|<t$, contradicting minimality. Thus $N_0\ne-1$.

For the second assertion put $y_0=V+e$. The exact identity
\[
 N_a=UV^2-(U-1)y_0^2
\]
is immediate. If $U=0$, then $W_s=I^2\geq81$. Otherwise $U\geq2$,
and $N_a=0$ would force the consecutive integers $U,U-1$ to be
squares, unless $V=y_0=0$. The latter contradicts $e>0$. Hence
$N_a\ne0$. Equality in the second block forces $W_s=0$ and
$N_a=\pm1$. Now $U=I^2$, so modulo four $N_a$ is congruent to
$y_0^2$ when $I$ is even and to $V^2$ when $I$ is odd. Thus $N_a=-1$
is impossible. This proves all assertions in~\eqref{eq:min-blocks}.

At a zero of~\eqref{eq:short54}, both minimum blocks attain their
minima and the remaining square vanishes. Hence
\[
 N_0=N_a=1,\qquad W_s=0,\qquad N_1N_kN_t=1.
\]
The last equation makes $N_1$ an integer unit. Since
$\Delta=(a+2)^2-1\equiv0$ or $3\pmod4$, square residues exclude
$N_1=-1$. Thus $N_1=1$ and $N_kN_t=1$. These are precisely the
equality conditions of~\eqref{eq:short46}; conversely they imply both
equations. This is a complete tuplewise equivalence, uniform in the
fixed ports with $\beta>0$.

The saving is transparent at the circuit level. Starting from $Q_{46}$,
remove the $N_1^2$ gate and its summation gate, insert one multiplication
of $N_1$ into the unit product, and change the final constant from three
to two. This saves exactly one binary operation.

\subsection{Squared-modulus completeness for the ordinary strong form}

The historical gap construction used
$V_{\mathrm{old}}=c(T_{\mathrm{old}}f-1)-Rf^2$.
The identity
\[
 V=f^2(cT-R)-c=V_{\mathrm{old}}\big|_{T_{\mathrm{old}}=fT}
\]
proves soundness immediately. Completeness requires reconstructing an
auxiliary Pell ordinate, rather than assuming divisibility of every
old quotient.

A halting computation supplies a canonical normalized witness.
Transfer it to ordinary coordinates by
$i_{\mathrm{ordinary}}=\Delta i_{\mathrm{normalized}}$.
The inherited decoding supplies $c=\psi_A(R)$,
$A=a+2$, $R\equiv3\pmod4$, and
\[
 I^2=\Delta(f^2-1),\qquad c\mid I,
 \qquad f^2\equiv1\pmod c,
 \qquad \gcd(c,f)=1.
\]
For $R=2r+1$ the Pell polynomials satisfy
\[
 \chi_t(R)=tQ_r(t^2),\quad
 Q_r(0)=(-1)^rR,\quad
 Q_r(1-A^2)=(-1)^r\psi_A(R).
\]
These follow by the elementary second-order Pell recurrences. Since $r$
is odd, take
\[
 V=\frac{\chi_I(R)}{I},\qquad y=\psi_I(R).
\]
Both are positive integers, $V>1$, and their norm is one. The displayed
polynomial identities give $V\equiv-c\pmod{f^2}$ and
$V\equiv-R\pmod c$. Consequently
\[
 T=\frac{V+c+Rf^2}{cf^2}>0
\]
is an integer and satisfies the required quotient formula. Finally
\[
 (I^2-1)(y^2-V^2)=V^2-1>0
\]
implies $e=y-V>0$. The other coordinates remain unchanged. This proves
existence equivalence to the authenticated ordinary source and hence,
using its inherited compiler and K215 calibration, the ordinary
mathematical implications formalizable in ZFC:
\begin{align*}
 \text{$Q_{54}$ has no positive zero}&\longleftrightarrow
 \operatorname{Con}(\mathrm{ZFC}),\\
 \text{$Q_{46}$ has no positive zero}&\longleftrightarrow
 \operatorname{Con}(\mathrm{ZFC}).
\end{align*}

\subsection{Exact degrees and bounded comparison}

The factor degrees are $22,18,23,22,7,2$ for
$N_0,N_1,N_a,W_s,N_k,N_t$, respectively. The main-norm cancellation is
the identity, with $Z=X+\gamma H$,
\[
 (ac+Z)^2-(a^2+H)c^2=Z^2+2acZ-Hc^2.
\]
Writing $Q=\beta J$, $k_0=\eta+\zeta$ and
$C_1=Q-F-\alpha$, the highest homogeneous forms are
\begin{align*}
 [Q_{54}]_{54}
 &=64\gamma^2h^2w^6s^8k_0^2Q^{30}(wC_1-zQ)^2,\\
 [Q_{46}]_{46}
 &=4w^4s^6k_0^2T^2f^8e^2Q^{22}.
\end{align*}
Their respective coefficients of
$J^{32}\gamma^2h^2w^6s^8\eta^2z^2$ and
$J^{22}w^4s^6\eta^2T^2f^8e^2$ are
$64\beta^{32}$ and $4\beta^{22}$, strictly positive for the supplied
tuple. The exact symbolic audit verifies these full leading forms,
not merely a numerical specialization.

These points improve the local degree/operation frontier between the
reviewed $72/80$ and $75/44$ candidates. They do not improve the existing
66-operation record, the degree-44 point, or the smallest-variable
constructions. The proof files and source receipts make this a bounded
repository comparison, without a claim of worldwide priority or global
minimality.


\section{The authoritative 31-coordinate source}\label{app:source}

This appendix fixes the source completely and avoids a possible
transport-coordinate ambiguity. Here the names are local source
coordinates, separate from the nine inputs in the main construction.
Their order is
\begin{align*}
 &(J,F,\alpha,z,f,h,i,o,j,s,w,\tau,\eta,\zeta,y,\gamma,\\
 &\hspace{2em}N,E,P,c,a,X,D,C_2,L,F_2,V,g,I,B_{\rm aux},\Delta).
\end{align*}
All are positive in the source. Define the affine abbreviations
$q=\beta J+1$, $k=\eta+\zeta$, $H_s=4a+3$, and $C_s=q-F-\alpha$.
The 22 equations, in their authoritative order, are
\begin{align*}
 L&=q^2-1,&N&=q(L+1),\\
 X&=wq,&a&=E+sN,\\
 E&=X(a-E),&c&=k(a-E)+\eta,\\
 P&=E(c-\eta)+k,&\tau^2&=P(P-k)+1,\\
 D&=X+ac+\gamma H_s,&\Delta&=a^2+H_s,\\
 C_2&=c^2,&D^2&=\Delta C_2+1.
\end{align*}
The next three rows are
\begin{align*}
 k&=jc-V+1+hE,\\
 jc-V&=(q-F)(N-q)-C_sL+(M_C+qM_F)J,\\
 (K_0+w)C_s+(q-F)-z(q-1)&=1.
\end{align*}
The final seven rows are
\begin{align*}
 F_2&=f^2,&V&=of-c,\\
 g&=y^2-V^2,&I&=iC_2,\\
 B_{\rm aux}&=I^2,&B_{\rm aux}&=\Delta(F_2-1),\\
 y^2-B_{\rm aux}g&=1.
\end{align*}
Take left minus right, and replace each of these 31 coordinates by
$x_i+1$, to obtain $p_1,\ldots,p_{22}$. Each residual has degree at most
two in $\mathbf x$ after the fixed ports are specialized. In particular,
the original positive transport coordinate is used in row 15. The
shifted transport equation $(K_0+w)C_s=F+z\beta J$ belongs to a related
source but is not substituted into this definition.

The archive retains the original JSON and its historical metadata
unchanged. The active manifest explicitly selects the whole package-31
tuple; old metadata in preserved evidence files does not override it.
The shared evaluator's transport adjustment is checked symbolically
against this exact source.

\section{Provenance and an incorporation checklist}

This note was derived from the placed collection and incoming submissions
at the pin given in Section 1. The principal immediate sources are:
\begin{itemize}
\item \texttt{ZFC\_Diophantine\_Selective\_Quadratic\_2026-10-23.zip}:
the exact original 31-variable source, selective weights and targets,
sparse target padding, and parity-calibrated threshold.
\item \texttt{con-zfc-sparse-pell-124c07c0b7.zip}:
the positivity-first radical merger, CRT absorption, and
the majorant $k_*(8D^2-4D-1)$.
\item \texttt{ConZFC\_Sparse\_Support\_2026-10-23.zip}:
the ordinary gap minimum argument and its coherent coefficient binding.
\item \texttt{con\_zfc\_sparse\_arithmetic\_2026-10-23.zip}:
the squared-modulus Pell reconstruction used in Appendix A.
\end{itemize}
The source receipts, copied notices, and licenses are retained. Newly
authored material is marked separately; no claim is made to relicense
inherited files. The independent reviews were performed within this
AI-assisted session and are not external human endorsements.

For incorporation, the mathematical milestones can be separated cleanly:
first prove the abstract shifted-norm/free-coordinate lemma; then the
two-block valuation formula and threshold identity; then certify the
literal shared evaluator against the sparse source; finally connect that
construction to the existing compiler endpoint. Each step has a precise
statement and finite data boundary in the archive. The present note does
not modify the repository or claim that those future kernel proofs have
already been compiled.

\begin{thebibliography}{9}
\bibitem{Jones1982}
James P. Jones, \emph{Universal Diophantine equation},
Journal of Symbolic Logic \textbf{47} (1982), no. 3, 549--571.
\url{https://doi.org/10.2307/2273588}.

\bibitem{Repo}
Vladimir Reshetnikov et al., \emph{ProveIt}, snapshot
\texttt{4b4349064a750ba6347293b55874a7f92abdd0db}, accessed 23 October 2026.
\href{https://github.com/VladimirReshetnikov/ProveIt/tree/4b4349064a750ba6347293b55874a7f92abdd0db/Computability/DiophantineConsistency}{Pinned repository source}.

\bibitem{Incoming}
ProveIt incoming research archives, same snapshot.
\href{https://github.com/VladimirReshetnikov/ProveIt/tree/4b4349064a750ba6347293b55874a7f92abdd0db/docs/incoming}{Pinned incoming collection}.

\bibitem{RatioLean}
ProveIt, \emph{NineRatio.lean}, numerical Jones ratio interfaces, same snapshot.
\href{https://github.com/VladimirReshetnikov/ProveIt/blob/4b4349064a750ba6347293b55874a7f92abdd0db/Computability/HilbertTenthProblem/Lean/Diophantine/Paper1982/NineRatio.lean}{Pinned numerical theorem source}.

\bibitem{ExplicitZFC}
ProveIt, \emph{A closed executable equation for ZFC consistency},
\texttt{EXPLICIT\_EQUATION.md}, same snapshot. The separate checked
construction and its object-language formalization boundary.
\href{https://github.com/VladimirReshetnikov/ProveIt/blob/4b4349064a750ba6347293b55874a7f92abdd0db/Computability/DiophantineConsistency/EXPLICIT_EQUATION.md}{Pinned endpoint and proof-status document}.
\end{thebibliography}
\end{document}
